% !TEX root = ../aa_SoM_Chapters.tex

%% HARRIS: Chapter 10

% ö = \%c3\%b6
% ä = \%C3\%A4

\xdef\sectiontitle{\IIsectiontitle} %aiheuttama

%%%%%%%%%%%%%%%
\begin{frame}[label=2_7_a]%[label=2_6_TOC1]
\frametitle{\texorpdfstring{\culine[\sectiontitleulinecolor]{\sectiontitle}}{\sectiontitle} }

%\vspace{0.5cm}
%\begin{itemize}[leftmargin=0.5cm,itemsep=3mm]
%    \item[2.1] \hyperlink{2_1_a}{\IIaSubsectiontitle}
%    \item[2.2] \hyperlink{2_2_a}{\IIbSubsectiontitle}	
%    \item[2.3] \hyperlink{2_3_a}{\IIcSubsectiontitle}	
%    \item[2.4] \hyperlink{2_4_a}{\IIdSubsectiontitle}
%    \item[2.5] \hyperlink{2_5_a}{\IIeSubsectiontitle}
%    \item[2.6] \hyperlink{2_6_a}{\CDAlert[red]{\IIfSubsectiontitle}}
%    \item[2.7] \hyperlink{2_7_a}{\IIgSubsectiontitle}
%    \item[2.8] \hyperlink{2_8_a}{\IIhSubsectiontitle}
%    \item[2.9] \hyperlink{2_9_a}{\IIiSubsectiontitle}
%\end{itemize}

\vspace{0.2cm}
\begin{itemize}[leftmargin=0.5cm,itemsep=3mm]
    \item[2.1] \hyperlink{2_1_a}{\IIaaSubsectiontitle}
    \item[2.2] \hyperlink{2_2_a}{\IIaSubsectiontitle}
    \item[2.3] \hyperlink{2_3_a}{\IIbSubsectiontitle}	
    \item[2.4] \hyperlink{2_4_a}{\IIcSubsectiontitle}	
    \item[2.5] \hyperlink{2_5_a}{\IIdSubsectiontitle}
    \item[2.6] \hyperlink{2_6_a}{\IIeSubsectiontitle}
    \item[2.7] \hyperlink{2_7_a}{\CDAlert[red]{\IIfSubsectiontitle}}
    \item[2.8] \hyperlink{2_8_a}{\IIgSubsectiontitle}
    \item[2.9] \hyperlink{2_9_a}{\IIhSubsectiontitle}
    \item[2.10] \hyperlink{2_10_a}{\IIiSubsectiontitle}
\end{itemize}

\endofslide \end{frame}
%%%%%%%%%%%%%%%

\xdef\sectiontitle{\IIfSubsectiontitle} 

%%%%%%%%%%%%%%%
\begin{frame}
\frametitle{\texorpdfstring{\culine[\sectiontitleulinecolor]{\sectiontitle}}{\sectiontitle} }

\begin{itemize}[itemsep=1.8mm]
	\item \CDAlert[red]{Switching on/off} \appear{and \CDAlert[blue]{amplification}} \appear{are essential properties of the components in modern electronic circuits} \appear{(in electronic communication}\appear{, mobiles}\appear{, computers}\appear{, displays}\appear{, lighting, digital communication}\appear{, etc.)} 
	
	\item \Alert[fuchsia]{Simple resistors, capacitors and inductors are not able to do those things}\appear{, we need components based on semiconductor materials} 
	
	\item Such components include:
	\item[] \begin{itemize} \seminormal
	\item  \CDAlert[fuchsia]{Diode}
\item \CDAlert[fuchsia]{Transistor}
\item Metal–oxide–semiconductor \appear{field-effect transistor} \appear{(MOSFET)}
\item Light-emitting diode \appear{(LED)}
\item Detectors \appear{(incident light/electrons generate an electrical current)}
\item Charge-coupled devices \appear{(CCDs)}
\item \CDAlert[fuchsia]{Complementary metal–oxide–semiconductor (CMOS) technology}
\item $\ldots$
\end{itemize}


\end{itemize}


\endofslide 
%\endofsection
\end{frame}
%%%%%%%%%%%%%%%

%%%%%%%%%%%%%%%
\begin{frame}
\frametitle{\texorpdfstring{\culine[\sectiontitleulinecolor]{\sectiontitle}}{\sectiontitle} }
\framesubtitle{Hall effect}

\begin{itemize}
	\item In semiconductors charge carriers are of both polarities: \appear{negative \CDAlert[red]{electrons}} \appear{and positive \CDAlert[red]{holes}} 
	
	\item This can be experimentally verified by so called \CDAlert[red]{Hall effect}. \appear{A current is arranged through a doped semiconductor}\appear{, which is arranged into a magnetic field}\appear{. Lorentz force creates drifting of the charge carriers sideways}\appear{, perpendicular to the direction of current: }
\[
\vec{F}  \equal  \appear{q} \appear{(\vec{v}} \appear{\times \vec{B})}
\]

\end{itemize}

\xdef\loc{south}
\begin{tikzpicture}[remember picture,overlay] % anchor=south
    \node (A) [yshift=-0.25*\figYshift,xshift=0*\figXshift, anchor=\loc] at (current page.\loc) {\includegraphics[page=1,height=3.5cm]{Figures_booklet/SOM_Tipler_Solid_state_physics_figures/SOM_Tipler_SSP_figure_39_cropped.png}}; %scale=\figscale. % 8cm max height
\end{tikzpicture}

\endofslide 
%\endofsection
\end{frame}
%%%%%%%%%%%%%%%

%%%%%%%%%%%%%%%
\begin{frame}
\frametitle{\texorpdfstring{\culine[\sectiontitleulinecolor]{\sectiontitle}}{\sectiontitle} }
\framesubtitle{Hall effect}

\begin{itemize}
\item Since the polarity of the charge carrier in the formula is\appear{, in fact, taken into account twice} \appear{(sign of charge and direction of movement, sign of $v_d$)}\appear{, in both cases, the charge carriers will eventually drift into the same direction}\appear{, revealing their sign by the induced perpendicular electric field $E_{Hall}$}

\item This induced \CDAlert[fuchsia]{Hall field $E_{Hall}$} then results in a measurable potential difference of form 
\[
V_{\text {Hall}}  \equal  \appear{v_{d}} \appear{B} \appear{w} \appear{\,, \quad ~where} \qquad w=width~of~the~strip
\]
\end{itemize}

\xdef\loc{south}
\begin{tikzpicture}[remember picture,overlay] % anchor=south
    \node (A) [yshift=-0.25*\figYshift,xshift=0*\figXshift, anchor=\loc] at (current page.\loc) {\includegraphics[page=1,height=3.5cm]{Figures_booklet/SOM_Tipler_Solid_state_physics_figures/SOM_Tipler_SSP_figure_39_cropped.png}}; %scale=\figscale. % 8cm max height
\end{tikzpicture}

\endofslide 
%\endofsection
\end{frame}
%%%%%%%%%%%%%%%

%%%%%%%%%%%%%%%
\begin{frame}
\frametitle{\texorpdfstring{\culine[\sectiontitleulinecolor]{\sectiontitle}}{\sectiontitle} }
\framesubtitle{Holes as charge carriers}

\begin{itemize}
	\item As charge carriers, the concept of holes is more complicated than just current going to the direction opposite to the movement of electrons
	
	\item At zero temperature in a semiconducting solid\appear{, the valence band is full and the conduction band is empty} 
	\implies{ The material will not respond to \\
			an external electric field }
\end{itemize}
	
% 6.5 cm = midway, 10 cm = 3/4 
\vspace{2.5mm}
\begin{minipage}{7.8cm}
\begin{itemize}
	\item At non-zero $T$\appear{, with thermal excitations taking place in the solid}\appear{, an electron will move from valence band to conduction band}\appear{, leaving behind a hole} 
	\implies{ \Alert[red]{Electron} is then free to move within \\
			 the conduction band }
	\implies{ A \Alert[blue]{hole} is also free to move within \\
			 the valence band }
\end{itemize}
\end{minipage}
\vspace{2.5mm}

\xdef\loc{south east}
\begin{tikzpicture}[remember picture,overlay] % anchor=south
    \node (A) [yshift=-0.25*\figYshift,xshift=\figXshift, anchor=\loc] at (current page.\loc) {\includegraphics[page=1,height=6cm]{Figures_booklet/SOM_10_Bonding_figures/SOM_Bonding_Figure_35.png}}; %scale=\figscale. % 8cm max height
\end{tikzpicture}

\endofslide 
%\endofsection
\end{frame}
%%%%%%%%%%%%%%%


%%%%%%%%%%%%%%%
\begin{frame}
\frametitle{\texorpdfstring{\culine[\sectiontitleulinecolor]{\sectiontitle}}{\sectiontitle} }
\framesubtitle{Holes as charge carriers}

\begin{itemize}
	\item As charge carriers, the concept of holes is more complicated than just current going to the direction opposite to the movement of electrons
	
	\item At zero temperature in a semiconducting solid, the valence band is full and the conduction band is empty 
	\implies{ The material will not respond to \\
			an external electric field }
\end{itemize}
	
% 6.5 cm = midway, 10 cm = 3/4 
\vspace{2.5mm}
\begin{minipage}{7.8cm}
\begin{itemize}
	\item The hole will not 'actually' move\appear{, but it gives rise to certain freedom of electrons to move} 
	\item \Alert[fuchsia]{Another electron can move to the empty state}\appear{, and another one to replace it }
	\implies{ While the valence electrons themselves are not free\appear{, the unoccupied state, the hole}\appear{, is free to move} }
\end{itemize}
\end{minipage}
\vspace{2.5mm}

\xdef\loc{south east}
\begin{tikzpicture}[remember picture,overlay] % anchor=south
    \node (A) [yshift=-0.25*\figYshift,xshift=\figXshift, anchor=\loc] at (current page.\loc) {\includegraphics[page=1,height=6cm]{Figures_booklet/SOM_10_Bonding_figures/SOM_Bonding_Figure_35.png}}; %scale=\figscale. % 8cm max height
\end{tikzpicture}

\endofslide 
%\endofsection
\end{frame}
%%%%%%%%%%%%%%%

%%%%%%%%%%%%%%%
\begin{frame}
\frametitle{\texorpdfstring{\culine[\sectiontitleulinecolor]{\sectiontitle}}{\sectiontitle} }
\framesubtitle{Holes as charge carriers}

\begin{itemize}
	\item Note the energy states in connection with the position of the hole in the energy scheme:
\item[] \begin{itemize}[itemsep=3mm]
	 \item When the hole gets higher in the band\appear{, a lower state results}
	 \end{itemize}
\end{itemize}

% 6.5 cm = midway, 10 cm = 3/4 
\vspace{2.5mm}
\begin{minipage}{7.6cm}
\begin{itemize}
\item[] \begin{itemize}[itemsep=3mm]
	\item This phenomenon may be compared to a \CDAlert[fuchsia]{rising air bubble in water}. \appear{While the bubble gets up}\appear{, a part of water sinks down}\appear{, resulting in a lower energy for the system.}
\item As an electron tends to sink\appear{, the hole tends to "float".} \appear{In moving among the states in the valence band}\appear{, a hole's energy varies in the opposite sense form that of an electron}
\end{itemize}
\end{itemize}
\end{minipage}
\vspace{2.5mm}


\xdef\loc{south east}
\begin{tikzpicture}[remember picture,overlay] % anchor=south
    \node (A) [yshift=-0.25*\figYshift,xshift=\figXshift, anchor=\loc] at (current page.\loc) {\includegraphics[page=1,height=6cm]{Figures_booklet/SOM_10_Bonding_figures/SOM_Bonding_Figure_36.png}}; %scale=\figscale. % 8cm max height
\end{tikzpicture}


\endofslide 
%\endofsection
\end{frame}
%%%%%%%%%%%%%%%

%%%%%%%%%%%%%%%%
%\begin{frame}
%\frametitle{\texorpdfstring{\culine[\sectiontitleulinecolor]{\sectiontitle}}{\sectiontitle} }
%\framesubtitle{Note on the density of states}
%
%\begin{itemize}
%	\item We saw that the density of states for free electrons in a quantum well formed a parabolic curve: $D=D(E) \propto \sqrt{E}$, increasing with increasing energy. The actual shape of $D$, depends on structure of the crystal lattice and the position of the band. In fact, it needs to be that the density is almost symmetric for the band: 
%	\item Suppose that a band is completelty filled. If one electron is removed (perhaps to the conduction band...), then it is possible to say that a hole has been created, since a vacant state is there. When an external force is applied, some elecron may be moved into the vacant state. But by doing so, the electron leaves a new hole, corresponding to a state that was previously occupied. We can say that the hole moves in a manner opposite to electron, i.e. it 'floats'. The zero energy for the holes is at the top of the band and their energy is measured downwards, i.e. $\quad\left(E-E_{\min }\right) \rightarrow\left(E_{\max }-E\right)$
%	
%\end{itemize}
%
%%\xdef\loc{east}
%%\begin{tikzpicture}[remember picture,overlay] % anchor=south
%%    \node (A) [yshift=\figYshift,xshift=\figXshift, anchor=\loc] at (current page.\loc) {\includegraphics[page=1,height=7cm]{Figures_booklet/SOM_10_Bonding_figures/SOM_Bonding_Figure_1.png}}; %scale=\figscale. % 8cm max height
%%\end{tikzpicture}
%
%\endofslide 
%%\endofsection
%\end{frame}
%%%%%%%%%%%%%%%%

%%%%%%%%%%%%%%%
\begin{frame}
\frametitle{\texorpdfstring{\culine[\sectiontitleulinecolor]{\sectiontitle}}{\sectiontitle} }
\framesubtitle{Effective mass of the charge carrier}

\seminormal
\begin{itemize}
	\item If a moving particle is treated as a wave\appear{, its velocity is described using so-called group velocity:}\vspace{-4mm}
\[
v_{\text {particle }}  \equal \appear{v_{\text {group }}}  \equal \frac{d \omega}{d k}
\]
\item The functional form of $\omega=\omega(k)$ in each case\appear{, and in each medium}\appear{, is called "\CDAlert[fuchsia]{dispersion relation}".} \appear{The acceleration of a particle would be:}
\[
a  \equal \frac{d \textrm{((}v_{\text {hiukkanen}} \textrm{))}}{d t}  \equal \frac{d}{d t} \LP \frac{d \omega}{d k} \,\RP  \equal \frac{d^{2} \omega}{d k^{2}} \frac{d k}{d t}
\]

\item Consider an external force acting on the particle\appear{. Then the work done by the force per unit time (power) is given by force times velocity:}
$$
\begin{aligned}
F_{\text {ext}} \appear{v_{\text {particle}}} & \equal \frac{d E}{d t}  \equal \frac{d(\hbar \omega)}{d t}  \equal \appear{\hbar} \frac{d \omega}{d t} \equal \appear{\hbar} \frac{d \omega}{d k} \frac{d k}{d t}  \equal \appear{\hbar v_{\text {particle}}} \frac{d k}{d t} \qquad \\
\Rightarrow \qquad \qquad \appear{F_{ext}}&  \equal \appear{\hbar} \frac{d k}{d t}
\end{aligned}
$$
\end{itemize}

\endofslide 
%\endofsection
\end{frame}
%%%%%%%%%%%%%%%

%%%%%%%%%%%%%%%
\begin{frame}
\frametitle{\texorpdfstring{\culine[\sectiontitleulinecolor]{\sectiontitle}}{\sectiontitle} }
\framesubtitle{Effective mass of the charge carrier}

\begin{itemize}
	\item We now use the form of Newton's 2nd law \appear{$(F=m a)$} \appear{to define \CDAlert[fuchsia]{effective mass} of the charge carrier:}
\end{itemize}
% 6.5 cm = midway, 10 cm = 3/4 
%\vspace{2.5mm}
\begin{minipage}{8.2cm}
\begin{itemize}
	\item[] 
	$$
	\begin{aligned}
m^{*} & \equal \frac{F_{e x t}}{a}  \equal \frac{\hbar \frac{d k}{d t}}{\frac{d^{2} \omega}{d k^{2}} \frac{d k}{d t}}  \equal \frac{\hbar}{\frac{d^{2} \omega}{d k^{2}}}  \equal \hbar^{2} \appear{\textrm{((}} \frac{d^{2} E}{d k^{2}} \appear{\textrm{))}^{-1}} \qquad \\
\Rightarrow \qquad \frac{1}{m^{*}}&  \equal \frac{1}{\hbar^{2}} \frac{d^{2} E}{d k^{2}}
\end{aligned}
$$

\item Let's verify that this choice for the definition for mass complies with the energy of the (free) particle:
\[
 E  \equal \frac{p^{2}}{2 m_e}  \equal \frac{\hbar^{2} k^{2}}{2 m_e}
\]
\end{itemize}
\end{minipage}
\vspace{2.5mm}


\xdef\loc{south east}
\begin{tikzpicture}[remember picture,overlay] % anchor=south
    \node (A) [yshift=-0.25*\figYshift,xshift=\figXshift, anchor=\loc] at (current page.\loc) {\includegraphics[page=1,height=7cm]{Figures_booklet/SOM_10_Bonding_figures/SOM_Bonding_Figure_27.png}}; %scale=\figscale. % 8cm max height
\end{tikzpicture}

\endofslide 
%\endofsection
\end{frame}
%%%%%%%%%%%%%%%


%%%%%%%%%%%%%%%
\begin{frame}
\frametitle{\texorpdfstring{\culine[\sectiontitleulinecolor]{\sectiontitle}}{\sectiontitle} }
\framesubtitle{Effective mass of the charge carrier}

% 6.5 cm = midway, 10 cm = 3/4 
%\vspace{2.5mm}
\begin{minipage}{8.2cm}
\begin{itemize}
	\item[] %\vspace{-4mm}
	$$
 E =\frac{p^{2}}{2 m_e}=\frac{\hbar^{2} k^{2}}{2 m_e}
$$
\item Let's look at Figure 27\appear{, showing the energy of the electrons \textrm{((}mass $m_{e} $\textrm{))} in a solid as a function of wavenumber $k$} 

\item The second derivative of energy with respect to $k$, will then also give:
\[
\frac{d^{2} E}{d k^{2}}  \equal  \appear{\hbar^{2}} \frac{1}{m_{e}} \quad \Rightarrow \quad \appear{m_{e}} \equal \appear{\hbar^{2}} \appear{\textrm{((}} \frac{d^{2} E}{d k^{2}} \appear{\textrm{))}^{-1}} \equiv \appear{m^*}
\]
\item Therefore the concept of an effective mass $m^*$ \appear{agrees with the mass for an electron}\appear{, for the case when the electron behaves like a free particle} \appear{(in the middle of the bands)}
\end{itemize}
\end{minipage}
\vspace{2.5mm}


\xdef\loc{south east}
\begin{tikzpicture}[remember picture,overlay] % anchor=south
    \node (A) [yshift=-0.25*\figYshift,xshift=\figXshift, anchor=\loc] at (current page.\loc) {\includegraphics[page=1,height=7cm]{Figures_booklet/SOM_10_Bonding_figures/SOM_Bonding_Figure_27.png}}; %scale=\figscale. % 8cm max height
\end{tikzpicture}

\endofslide 
%\endofsection
\end{frame}
%%%%%%%%%%%%%%%

%%%%%%%%%%%%%%%
\begin{frame}
\frametitle{\texorpdfstring{\culine[\sectiontitleulinecolor]{\sectiontitle}}{\sectiontitle} }
\framesubtitle{Effective mass of the charge carrier}

\semismall
\begin{itemize}
%	\item[] \vspace{-2mm}
%	$$
%m_{e}=\hbar^{2} \textrm{((}\Frac{d^{2} E}{d k^{2}} \textrm{))}^{-1} \qquad \& \qquad \frac{1}{m_{e}} \equal \frac{1}{\hbar^{2}} \frac{d^{2} E}{d k^{2}}
%$$
\item The curvature \appear{(= second derivative)} \appear{of $E=E(k)$} \appear{gives the reciprocal mass (inverse of mass) }
\end{itemize}

\vspace{1.5mm}

\begin{minipage}{7.8cm}
\begin{itemize}[itemsep=1.5mm]

\item When the pink curve follows the dashed curve\appear{, i.e. the parabola $E \propto k^{2}$}\appear{, then the energy follows that of a free electron}\appear{, and we have $m^{*}=m_{e}$}

\item Near band edges\appear{, on both sides of the gap}\appear{, effective mass deviates from $m_{e}$.} \appear{Starting from $k=0$, near $k=\pi / a$ the curvature becomes large} \appear{and very close to the boundary even negative.} \appear{Just above the gap the curvature is large and positive. And so on.}

\item For the regions, where the gap energy is small compared to the width of the band\appear{, \Alert[blue]{effective masses $m^{*}$ may become $0.01$ or $0.1$ times the mass of a free electron}}
\end{itemize}
\end{minipage}

\xdef\loc{south east}
\begin{tikzpicture}[remember picture,overlay] % anchor=south
    \node (A) [yshift=-0.25*\figYshift,xshift=\figXshift, anchor=\loc] at (current page.\loc) {\includegraphics[page=1,height=7cm]{Figures_booklet/SOM_10_Bonding_figures/SOM_Bonding_Figure_37.png}}; %scale=\figscale. % 8cm max height
\end{tikzpicture}


\endofslide 
%\endofsection
\end{frame}
%%%%%%%%%%%%%%%

%%%%%%%%%%%%%%%
\begin{frame}
\frametitle{\texorpdfstring{\culine[\sectiontitleulinecolor]{\sectiontitle}}{\sectiontitle} }
\framesubtitle{Conductivity due to both e's and h's}

\seminormal

\begin{itemize}
	\item In an ordinary conductor (such as copper)\appear{, electrons fill the states around the middle of the band}\appear{, where they are essentially free particles of mass $m$} 
	\item In a semiconductor\appear{, electrons fill the states to the top of valence band}\appear{, where the effective mass is negative }
	\item Hole energy varies in the opposite sense\appear{ (because the valence band is oppositely curved than the conduction band),}
\end{itemize}

% 6.5 cm = midway, 10 cm = 3/4 
%\vspace{2.5mm}
\begin{minipage}{8.5cm}
\begin{itemize}
	\item[] therefore at the top of a band and only there, a hole has positive mass and thus\appear{, it moves in the direction of the external field}

\item In total, both conduction electrons and valence holes act as \CDAlert[fuchsia]{charge carriers} \appear{and respond to external field}\appear{, and contribute to electrical conductivity:}\vspace{-2.5mm}
\[
\sigma  \equal \frac{n_{e} e^{2} \tau_{e}}{m_{e}^{*}}  \plus \frac{n_{h} e^{2} \tau_{h}}{m_{h}^{*}}
\]
\end{itemize}
\end{minipage}

\xdef\loc{south east}
\begin{tikzpicture}[remember picture,overlay] % anchor=south
    \node (A) [yshift=\figYshift,xshift=\figXshift, anchor=\loc] at (current page.\loc) {\includegraphics[page=1,height=6cm]{Figures_booklet/SOM_10_Bonding_figures/SOM_Bonding_Figure_38.png}}; %scale=\figscale. % 8cm max height
\end{tikzpicture}

\endofslide 
%\endofsection
\end{frame}
%%%%%%%%%%%%%%%

%%%%%%%%%%%%%%%
\begin{frame}
\frametitle{\texorpdfstring{\culine[\sectiontitleulinecolor]{\sectiontitle}}{\sectiontitle} }
\framesubtitle{Number densities of charge carriers}

%
\begin{itemize}
	\item Recall the density of states $D=D(E)$ and occupation number $ \mathbb{N} = \mathbb{N} (E)_{F D}$ for fermions:
\[
D(E)  \equal \frac{(2 s+1) m^{3 / 2} V}{\pi^{2} \hbar^{3} \appear{\sqrt{2}}} \appear{\sqrt{E} }\qquad \appear{\&} \qquad  \appear{\mathbb{N} (E)_{F D}}  \equal \frac{1}{e^{ \textrm{((}E-E_{F} \textrm{))} / k_{B} T}+1}
\]

\item We can estimate the number density of the charge carriers in a semiconductor \appear{by integrating $D(E) \times  \mathbb{N} (E)_{F D}$ over an appropriate energy range}
\end{itemize}

% 6.5 cm = midway, 10 cm = 3/4 
\vspace{2.5mm}
\begin{minipage}{7cm}
\begin{itemize}
	\item Let $E_{c}$ and $E_{v}$ be the energies of the bottom of the conduction band and top of the valence band, respectively

\item In this case, the band gap energy is simply $\appear{E_g}  \equal \appear{E_{c}} \minus \appear{E_{v}}$
\end{itemize}
\end{minipage}

\xdef\loc{south east}
\begin{tikzpicture}[remember picture,overlay] % anchor=south
    \node (A) [yshift=-0.25*\figYshift,xshift=1*\figXshift, anchor=\loc] at (current page.\loc) {\includegraphics[page=1,height=4cm]{Figures_booklet/SOM_10_Bonding_figures/Tikz_SOM_Bonding_figures5}}; %scale=\figscale. % 8cm max height
\end{tikzpicture}

\endofslide 
%\endofsection
\end{frame}
%%%%%%%%%%%%%%%

%%%%%%%%%%%%%%%
\begin{frame}
\frametitle{\texorpdfstring{\culine[\sectiontitleulinecolor]{\sectiontitle}}{\sectiontitle} }
\framesubtitle{Number densities of charge carriers}

\begin{itemize}
	\item For \CDAlert[blue]{negative charge carriers} $(s=1 / 2)$ we write $D(E) \times  \mathbb{N} (E)_{F D}$:
\[
D(E) \appear{\times  \mathbb{N} (E)_{F D}} \equal \fraca{2 m_{e}^{* 3 / 2} V}{\pi^{2} \hbar^{3} \sqrt{2} \appear{\sqrt{E-E_{c}}}} \fraca{1}{e^{ \textrm{((}E-E_{F} \textrm{))} / k_{B} T}+1}
\]
\appear{We can use this directly in the following}

\item But, for \CDAlert[red]{positive charge carriers} $(s=1 / 2)$\appear{, i.e. the holes}\appear{, the shape of the band is different at the top of the valence band}\appear{, and so the square root of the energy needs to be replaced by $\sqrt{E_{v}-E}$}

\item Finally, the Fermi–Dirac distribution $  \mathbb{N} (E)_{F D}=\Frac{1}{e^{ \textrm{((}E-E_{F} \textrm{))} / k_{B} T}+1}$
\appear{needs to be replaced by:} $\appear{1}  \minus  \frac{1}{e^{ \textrm{((}E-E_{F} \textrm{))} / k_{B} T}+1}$

\item The combined probability of hole occurring for certain energy\appear{, added to the probability of the state being occupied}\appear{, must be unity} \appear{(\Alert[fuchsia]{the electron must be \CDAlert[fuchsia]{somewhere}})}
%\item The probability of hole occurring for certain energy when added to the probability of the state being occupied, must be unity
\end{itemize}

\endofslide 
%\endofsection
\end{frame}
%%%%%%%%%%%%%%%

%%%%%%%%%%%%%%%
\begin{frame}
\frametitle{\texorpdfstring{\culine[\sectiontitleulinecolor]{\sectiontitle}}{\sectiontitle} }
\framesubtitle{Number densities of charge carriers}

\begin{itemize}
	\item We obtain:
\[
1 \minus \frac{1}{e^{\frac{E-E_{F}}{k_{B} T}}+1} \equal \frac{e^{\frac{E-E_{F}}{k_{B} T}}}{e^{\frac{E-E_{F}}{k_{B} T}}+1} \equal \frac{1}{1+e^{\frac{- \textrm{((}E-E_{F} \textrm{))}}{k_{B} T}}} \equal \frac{1}{1+e^{\frac{E_{F}-E}{k_{B} T}}} \approx \appear{e^{\frac{- \textrm{((}E_{F}-E \textrm{))}}{k_{B} T}}}
\]

\item So, for positive charge carriers:
\[
D(E) \appear{\times  \mathbb{N} (E)_{F D}} \equal \frac{2 m_{h}^{* 3 / 2} V}{\pi^{2} \hbar^{3} \sqrt{2}} \appear{\sqrt{E_{v}-E}} \appear{e^{- \textrm{((}E_{F}-E \textrm{))} /k_{B} T}}
\]
\item Now, we can calculate the number densities for negative and positive charge carriers
\end{itemize}

\endofslide 
%\endofsection
\end{frame}
%%%%%%%%%%%%%%%

%%%%%%%%%%%%%%%
\begin{frame}
\frametitle{\texorpdfstring{\culine[\sectiontitleulinecolor]{\sectiontitle}}{\sectiontitle} }
\framesubtitle{Number densities of charge carriers}

\begin{itemize}
	\item For negative charge carriers we get $n_{e}$:
$$
\begin{aligned} 
\appear{n_{e}}&  \equal  \appear{\nint_{E_{c}}^{\infty}} \frac{2 m_{e}^{* 3 / 2} V}{\pi^{2} \hbar^{3} \sqrt{2}} \appear{\sqrt{E-E_{c}}} \appear{e^{\frac{- \textrm{((}E-E_{F} \textrm{))}}{k_{B} T}}} \appear{d E}  \equal \frac{2 m_{e}^{* 3 / 2} V}{\pi^{2} \hbar^{3} \sqrt{2}} \appear{e^{E_{F / k_{B}}}} \appear{\nint_{E_{c}}^{\infty}} \appear{\sqrt{E-E_{c}}} \appear{e^{\frac{-E}{k_{B} T}} d E} \\ 
& \equal \frac{1}{4} \appear{\textrm{((}} \frac{2 m_{e}^{*} k_{B} T}{\pi \hbar^{2}} \appear{\textrm{))}} \appear{^{3 / 2}} \appear{e^{- \textrm{((}E_{c}-E_{F} \textrm{))} / k_{B} T} }
\end{aligned}
$$

\item And for positive charge carriers $n_{h}$:
$$
\begin{aligned} 
\appear{n_{h}}&  \equal  \appear{\nint_{-\infty}^{E_{v}}} \frac{2 m_{h}^{* 3 / 2} V}{\pi^{2} \hbar^{3} \sqrt{2}} \appear{\sqrt{E_{v}-E}} \appear{e^{\frac{- \textrm{((}E_{F}-E \textrm{))}}{k_{B} T}} d E}  \equal \frac{2 m_{e}^{* 3 / 2} V}{\pi^{2} \hbar^{3} \sqrt{2}} \appear{e^{-E_{F / k_{B}}} } \appear{\nint_{-\infty}^{E_{v}}} \appear{\sqrt{E_{v}-E}} \appear{e^{\frac{E}{k_{B} T}} d E  }\\ 
& \equal \frac{1}{4} \appear{\textrm{((}} \frac{2 m_{h}^{*} k_{B} T}{\pi \hbar^{2}} \appear{\textrm{))}^{3 / 2}} \appear{e^{- \textrm{((}E_{F}-E_{v} \textrm{))} / k_{B} T}} \\ 
\end{aligned}
$$
\item Verification of these results in the Exercises !! 
\end{itemize}

\endofslide 
%\endofsection
\end{frame}
%%%%%%%%%%%%%%%

%%%%%%%%%%%%%%%
\begin{frame}
\frametitle{\texorpdfstring{\culine[\sectiontitleulinecolor]{\sectiontitle}}{\sectiontitle} }
\framesubtitle{Mass action law: \quad $n_e n_h =n_i^2$}

\begin{itemize}
	\item So now, as an overall result we obtain:
\[
n_{e} \appear{n_{h}}  \equal \frac{1}{2} \appear{\textrm{((}} \frac{k_{B} T}{\pi \hbar^{2}} \appear{\textrm{))}^{3}} \appear{\textrm{((}m_{e}^{*} m_{h}^{*} \textrm{))}^{3 / 2}} \appear{e^{-\Delta E / k_{B} T}} \appear{\,, \qquad \Delta E}  \equal  \appear{E_g}  \equal  \appear{E_c - E_v}
\]

\item Note, for intrinsic semiconductors \appear{the number densities of negative and positive charge carriers are equal:}
\[
n_{e}  \equal  \appear{n_{h}}  \equal  \frac{1}{4} \appear{\textrm{((}} \frac{2 k_{B} T}{\pi \hbar^{2}} \appear{\textrm{))}^{3 / 2}} \appear{\textrm{((}m_{e}^{*} m_{h}^{*} \textrm{))}^{3 / 4}} \appear{e^{-\Delta E / 2 k_{B} T}} \equiv \appear{n_i^2}
\]
\item[] where $n_i$ is the carrier density of an intrinsic semiconductor\appear{, using which we can write the highly useful \CDAlert[red]{mass action law}:} 
\[
n_{e} \appear{n_{h}}  \equal \appear{n_i} \qquad \appear{\textrm{(applies also to \Alert[red]{extrinsic semiconductors}!)}}
\]
\item Note also, that charge carriers numbers vary with both temperature \appear{and energy gap $\Delta E$}\appear{, i.e. depend also on the material}
\end{itemize}

\endofslide 
%\endofsection
\end{frame}
%%%%%%%%%%%%%%%

%%%%%%%%%%%%%%%
\begin{frame}
\frametitle{\texorpdfstring{\culine[\sectiontitleulinecolor]{\sectiontitle}}{\sectiontitle} }

\framesubtitle{Intrinsic semiconductors}
% On intrinsic semiconductors

\seminormal
\begin{itemize}[itemsep=1.6mm]
%	\item Previously, we discussed materials which are defined as \Alert[fuchsia]{intrinsic semiconductors}\appear{, undoped elements} \appear{(or compounds)} \appear{with a band gap smaller than $2$\;eV} \appear{(the choice of $2$\;eV is somewhat arbitrary)} 
	
	\item We have so far discussed of \CDAlert[red]{elemental semiconductors} \appear{(e.g. Si}\appear{, Ge)}\appear{, or \CDAlert[blue]{compound semiconductors}} \appear{(e.g. SiC}\appear{, Si$_{1-x}$Ge$_{x}$}\appear{, GaAs)}\appear{, which have all been \CDAlert[fuchsia]{intrinsic semiconductors}}\appear{, i.e. undoped semiconductors}  \appear{with a band gap smaller than $2$\;eV} \appear{(the choice of $2$\;eV is somewhat arbitrary)} 
	
	\item With such a small gap\appear{, the temperature does not have to be very high to produce significant occupation of the conduction band}
	
	\end{itemize}
% 6.5 cm = midway, 10 cm = 3/4 
\vspace{2.5mm}
\begin{minipage}{5.8cm}
\begin{itemize}[itemsep=1.6mm]
%\item[] significant occupation of the conduction band
	\item Two very representative intrinsic semiconductors are \CDAlert[red]{silicon} \appear{and \CDAlert[blue]{germanium},}%with band gaps of $1.14$\;eV and $0.67$\;eV, respectively
	%\item[] 
	\appear{ with band gaps of \CDAlert[red]{$1.14$\;eV}} \appear{and \CDAlert[blue]{$0.67$\;eV}, respectively}

\item Both of them are "covalent" solids\appear{, where all the atoms share all the 4 valence electrons}\appear{, for silicon in the $\mathrm{n}=3$ shell}\appear{, for germanium $(Z=32)$} \appear{in the $n=4$ shell}
\end{itemize}
\end{minipage}
\vspace{2.5mm}


\xdef\loc{south east}
\begin{tikzpicture}[remember picture,overlay] % anchor=south
    \node (A) [yshift=-0.25*\figYshift,xshift=\figXshift, anchor=\loc] at (current page.\loc) {\includegraphics[page=1,height=5cm]{Figures_booklet/SOM_Tipler_Solid_state_physics_figures/SOM_Tipler_SSP_figure_36_cropped.png}}; %scale=\figscale. % 8cm max height
\end{tikzpicture}

\endofslide 
%\endofsection
\end{frame}
%%%%%%%%%%%%%%%

%%%%%%%%%%%%%%%
\begin{frame}
\frametitle{\texorpdfstring{\culine[\sectiontitleulinecolor]{\sectiontitle}}{\sectiontitle} }
\framesubtitle{Intrinsic semiconductors}

\seminormal
\begin{itemize}[itemsep=2.2mm]
	\item The covalent bond between two Si atoms with the high electron density near the two neighbouring atoms \appear{is shown by arrow in the X-ray scattering measurements in the previous Figure}

\item Arsenic \appear{$(As, Z=33)$} \appear{in amorphous form also qualified as an elemental semiconductor with a band gap of $1.2-1.4$\,eV.} \appear{However As has \CDAlert[red]{five valence electrons}} \appear{(i.e. is \CDAlert[red]{pentavalent})} \appear{in the $n=4$ shell.} \appear{Four of the five valence electrons take part in the covalent bonds}\appear{, but the \CDAlert[red]{fifth electron} is very \CDAlert[red]{loosely bound} to the atom}

\item As is more often used in making of \CDAlert[blue]{III-V semiconductors} (e.g. GaAs), or as a \CDAlert[red]{dopant element in making of extrinsic semiconductors} 

\item When As is used as a dopant, the fifth electron occupies an energy level that is just slightly below the conduction band\appear{, but it can be easily excited into the conduction band.} \appear{This fifth electron and the As ion core form a \CDAlert[tunired]{hydrogen-like system}}\appear{, with approximate energies close to it}

%\item In pure materials, the number of electrons thermally excited to the conduction band \appear{necessarily equals the number of holes left behind in the valence band}
\end{itemize}

\endofslide 
%\endofsection
\end{frame}
%%%%%%%%%%%%%%%

%%%%%%%%%%%%%%%
\begin{frame}
\frametitle{\texorpdfstring{\culine[\sectiontitleulinecolor]{\sectiontitle}}{\sectiontitle} }

\framesubtitle{Extrinsic semiconductors}
%Doping of semiconductors $-$ Extrinsic semiconductors

\begin{itemize}
	\item By doping\appear{, i.e. adding small amounts of impurity elements}\appear{, we produce \CDAlert[fuchsia]{extrinsic semiconductors}} 
	
	\item These have a \Alert[fuchsia]{large majority of either conduction electrons or valence holes} \appear{and, in either case}\appear{, much greater concentration of charge carriers than the pure intrinsic semiconductor}

\item As an example, we present covalent Si \appear{(or Ge lattice)}\appear{, \CDAlert[red]{doped with the pentavalent As}}
\end{itemize}

\appear{\xdef\loc{south}
\begin{tikzpicture}[remember picture,overlay] % anchor=south
    \node (A) [yshift=-0.25*\figYshift,xshift=0*\figXshift, anchor=\loc] at (current page.\loc) {\includegraphics[page=2,height=4cm]{Figures_booklet/SOM_10_Bonding_figures/Tikz_SOM_Bonding_figures6_12_fixed}}; %scale=\figscale. % 8cm max height
\end{tikzpicture}
}

\endofslide 
%\endofsection
\end{frame}
%%%%%%%%%%%%%%%

%%%%%%%%%%%%%%%
\begin{frame}
\frametitle{\texorpdfstring{\culine[\sectiontitleulinecolor]{\sectiontitle}}{\sectiontitle} }

\begin{itemize}
	\item Let us continue analyzing the fifth electron of As 
	
	\item \CDAlert[tunired]{Bohr theory} \appear{can be used to estimate the energies available to it:}
	\[
E_{n}  \equal  \minus \frac{1}{2} \appear{\textrm{((}} \frac{k e^{2}}{\hbar} \appear{\textrm{))}} \frac{m_{e}}{n^{2}} \appear{\times} \frac{m^{*}}{m_{e}}  \frac{1}{\kappa^{2}} \equal  \minus \frac{1}{2} \appear{\textrm{((}} \frac{k e^{2}}{\hbar} \appear{\textrm{))}} \frac{m^{*}}{\kappa^{2}}  \frac{1}{n^{2}}
\]
\appear{here $\mathrm{m}^{*}$ is the effective mass}\appear{, $\kappa$ is the dielectric constant}\appear{, and the average radii of the Bohr orbits are obtained from} 
\[
 < r_{n}>\,  \equal \appear{a_{0}} \appear{n^{2}}  \frac{m_{e}}{m^{*}}  \appear{\kappa} \appear{\,, \qquad a_{0}} \appear{= Bohr~radius~0.0529~\mathrm{nm} }
 \] 

\item This energy is relative to the minimum energy of the conduction band $E_C$\appear{, since the As atom is introduced as an impurity to the Si lattice }
\implies{ one extra electron is free to move and to conduct electricity
 }
\item $E_{\infty}=0$ is at the energy $E_C$\appear{, other energies $E_{n}$ lie close below it, inside the band gap}

\end{itemize}

\endofslide 
%\endofsection
\end{frame}
%%%%%%%%%%%%%%%

%%%%%%%%%%%%%%%
\begin{frame}
\frametitle{\texorpdfstring{\culine[\sectiontitleulinecolor]{\sectiontitle}}{\sectiontitle} }
\framesubtitle{n-type extrinsic semiconductor}

\seminormal
\begin{itemize}[itemsep=1.65mm]
	
\item If we assume $m^* \approx 0.2 m_{e}$ \appear{(value for $m^{*}$ in silicon)} \appear{then $E_{1}=-0.020$\;eV lies lower relative the bottom of conduction band} \appear{i.e. is 'smaller' than the $-13.6$\;eV (= hydrogen ground state)}

\item Also $\appear{<r_{1}>=3.1 \mathrm{~nm}} \appear{=60 \times}$ \appear{radius of hydrogen ground state}

\item Thus, the extra electrons can be easily excited to the conduction band\appear{, since the \CDAlert[red]{ionization energy is close to $k T$ of room temperature}}

\item The hydrogen-like levels just below the conduction band are called \CDAlert[red]{donor levels}\appear{, since they \Alert[red]{'donate' electrons to the conduction band}} \appear{without leaving holes in the valence band}

\item This type of doped semiconductor is called \CDAlert[red]{n-type semiconductor}

\item The conductivity of the doped material can be controlled by changing the impurity concentration
\implies{ Already few parts per million \CDAlert[red]{(ppm) of impurity} noticeably change the conductivity }
\end{itemize}

\endofslide 
%\endofsection
\end{frame}
%%%%%%%%%%%%%%%

%%%%%%%%%%%%%%%%
%\begin{frame}
%\frametitle{\texorpdfstring{\culine[\sectiontitleulinecolor]{\sectiontitle}}{\sectiontitle} }
%
%%n-tyoe extrinsic semiconductor
%\begin{itemize}
%	\item Thus, the extra electrons can be easily excited to the conduction band, since the ionization energy is close to $k T$ of room temperature
%
%\item The hydrogen-like levels just below the conduction band are called donor levels, since they 'donate' electrons to the conduction band without leaving holes in the valence band.
%
%\item This type of impurity semiconductor is called $\underline{n}$-type semiconductor
%
%\item The conductivity of the doped material can be controlled by controlling the amount of impurity added
%
%\item Even few parts per million (ppm) of impurity
%\end{itemize}
%
%%\xdef\loc{east}
%%\begin{tikzpicture}[remember picture,overlay] % anchor=south
%%    \node (A) [yshift=\figYshift,xshift=\figXshift, anchor=\loc] at (current page.\loc) {\includegraphics[page=1,height=7cm]{Figures_booklet/SOM_10_Bonding_figures/SOM_Bonding_Figure_1.png}}; %scale=\figscale. % 8cm max height
%%\end{tikzpicture}
%
%\endofslide 
%%\endofsection
%\end{frame}
%%%%%%%%%%%%%%%%

%%%%%%%%%%%%%%%
\begin{frame}
\frametitle{\texorpdfstring{\culine[\sectiontitleulinecolor]{\sectiontitle}}{\sectiontitle} }
\framesubtitle{n-type extrinsic semiconductor, summary}

\begin{itemize}
	\item The valence band of the intrinsic semiconductor is completely full \implies{ Each impurity atom add extra electron in the conduction band }

\item Also, each impurity atom has an extra positive charge

\item The net effect is 'donor state'\appear{, and it lies just below the conduction band}\appear{, typically around ${\sim} 0.05$\;eV lower in energy}
\end{itemize}
\vspace{2mm}
\appear{\xdef\loc{south east}%
\begin{tikzpicture}[remember picture,overlay]% anchor=south
    \node (A) [yshift=-0.25*\figYshift,xshift=\figXshift, anchor=\loc] at (current page.\loc) {\includegraphics[page=1,width=7cm]{Figures_booklet/SOM_10_Bonding_figures/SOM_Bonding_Figure_39.png}};%scale=\figscale. % 8cm max height
\end{tikzpicture}%
}

\begin{minipage}{6.5cm}
\begin{itemize} 
	\item Therefore, the electrons from donor states \appear{can be excited to the conduction band} \appear{already by fairly low ambient temperatures} 
\end{itemize}

\end{minipage}

\endofslide 
%\endofsection
\end{frame}
%%%%%%%%%%%%%%%

%%%%%%%%%%%%%%%
\begin{frame}
\frametitle{\texorpdfstring{\culine[\sectiontitleulinecolor]{\sectiontitle}}{\sectiontitle} }
\framesubtitle{p-type extrinsic semiconductor}

\seminormal
%p-type extrinsic semiconductor
\begin{itemize}[itemsep=1.6mm]
	\item Another type of an extrinsic semiconductor is a \CDAlert[blue]{$\mathbf{p}$-type semiconductor}
	
	\item There, a silicon atom in the lattice can be replaced by a \CDAlert[blue]{trivalent} \Alert[blue]{aluminum \appear{(or gallium)}} \appear{atom}\appear{, which has three valence electrons}
	
	\item Aluminum accepts electrons from the valence band of silicon to complete four covalent bonds
	\implies{  These are '\CDAlert[blue]{acceptor states}' }

	\item Acceptor states are easily filled by electrons from the valence band\appear{, excited already by $k_B T$ of room temperature}
\end{itemize}

\appear{\xdef\loc{south}
\begin{tikzpicture}[remember picture,overlay] % anchor=south
    \node (A) [yshift=-0.25*\figYshift,xshift=0*\figXshift, anchor=\loc] at (current page.\loc) {\includegraphics[page=3,height=4cm]{Figures_booklet/SOM_10_Bonding_figures/Tikz_SOM_Bonding_figures6_12_fixed}}; %scale=\figscale. % 8cm max height
\end{tikzpicture}
}

\endofslide 
%\endofsection
\end{frame}
%%%%%%%%%%%%%%%

%%%%%%%%%%%%%%%
\begin{frame}
\frametitle{\texorpdfstring{\culine[\sectiontitleulinecolor]{\sectiontitle}}{\sectiontitle} }
\framesubtitle{p-type extrinsic semiconductor}

\begin{itemize}
	\item When an electron is excited to the acceptor states\appear{, a hole is left behind into the valence band}
	\implies{ Positive charge carriers are created into the semiconductor }
	
	\item To summarize, depending on the type of the dopant\appear{, i.e. the added impurities}\appear{, an intrinsic semiconductor \Alert[fuchsia]{can be made either n- or p-type extrinsic semiconductor}}
\end{itemize}

% 6.5 cm = midway, 10 cm = 3/4 
\vspace{2.5mm}
\begin{minipage}{6.5cm}
	\begin{itemize}	
	\item The \Alert[red]{conductivities can be greatly increased}! 
	
	\item More importantly\appear{, sophisticated combinations of n- and p-type extrinsic semiconductors} \appear{can be used to make marvellous \Alert[fuchsia]{electronic components}! }
\end{itemize}
\end{minipage}

\xdef\loc{south east}
\begin{tikzpicture}[remember picture,overlay] % anchor=south
    \node (A) [yshift=-0.25*\figYshift,xshift=1*\figXshift, anchor=\loc] at (current page.\loc) {\includegraphics[page=1,width=7cm]{Figures_booklet/SOM_10_Bonding_figures/SOM_Bonding_Figure_40.png}}; %scale=\figscale. % 8cm max height
\end{tikzpicture}

%\endofslide 
\endofsection
\end{frame}
%%%%%%%%%%%%%%%

